% chapters/03/praktikum.tex
\begin{praktikumbox}
Praktikum ini mengikuti Tugas Praktikum 01 pada RPS, yaitu menyiapkan
lingkungan kerja, menjalankan program jendela pertama, dan menyimpan struktur
proyek ke repositori Git. Kerjakan seluruh langkah secara berurutan dan isi
lembar kerja pada Tabel~\ref{tab:lembar-praktikum-01}. Kumpulkan hasilnya pada
repositori Git pribadi Anda.

\begin{langkah}
  \item \textbf{Periksa lingkungan kerja.} Jalankan skrip
        \code{ch03_periksa_lingkungan.py} dari folder \code{code}. Catat
        seluruh keluarannya, termasuk baris \istilah{Penafsir Python} dan
        \istilah{Pengandar video}. Apabila ada bagian yang belum siap,
        selesaikan lebih dahulu dengan merujuk
        Tabel~\ref{tab:masalah-pemasangan}.

  \item \textbf{Buat struktur folder.} Buat folder proyek bernama
        \code{proyek-pygame} dengan susunan seperti pada
        Gambar~\ref{fig:struktur-proyek}: folder \code{code},
        \code{catatan}, dan \code{aset}.

  \item \textbf{Salin skrip pemeriksaan} ke dalam folder \code{code}, lalu
        pastikan skrip tersebut masih berjalan dari tempat yang baru.

  \item \textbf{Ketik ulang program jendela.} Ketik ulang
        \code{ch03_jendela.py} dari Daftar~\ref{lst:ch03-jendela} ke dalam
        folder \code{code}. \textbf{Jangan menyalin dan menempelkan} berkas
        yang tersedia; mengetik ulang membuat Anda memperhatikan setiap baris.
        Jalankan program, lalu tutup jendelanya.

  \item \textbf{Ubah ukuran dan warna.} Jalankan program dengan lebar 800 dan
        tinggi 600 piksel. Setelah itu, ubah nilai \code{WARNA_LATAR} menjadi
        warna pilihan Anda, lalu jalankan kembali.

  \item \textbf{Tambahkan satu baris teks} berisi nama dan nomor induk
        mahasiswa Anda, lalu atur posisinya agar tidak bertumpuk dengan baris
        lain.

  \item \textbf{Buat berkas \code{.gitignore}} sesuai
        Bagian~\ref{sec:struktur-git}, lalu simpan seluruh pekerjaan dengan
        pesan penyimpanan yang menerangkan isinya.

  \item \textbf{Tulis catatan praktikum} pada folder \code{catatan} memuat
        kendala yang Anda hadapi dan cara penyelesaiannya. Catatan ini menjadi
        bagian penilaian kelengkapan dokumentasi.
\end{langkah}
\end{praktikumbox}

\begin{table}[htbp]
\centering
\caption{Lembar kerja Praktikum 01}
\label{tab:lembar-praktikum-01}
{\setlength{\tabcolsep}{4pt}\footnotesize%
\begin{tabularx}{\textwidth}{@{}C{1.6cm} L{5.0cm} Y@{}}
\toprule
\headrow \textbf{Langkah} & \textbf{Hasil yang dicatat} & \textbf{Keterangan Anda} \\
\midrule
1 & Versi Python, versi Pygame, versi SDL, dan nama pengandar video & \\[1.2em]
\addlinespace
1 & Nama berkas penafsir Python yang dipakai & \\[1.2em]
\addlinespace
4 & Apakah jendela terbuka dan tertutup dengan benar? & \\[1.2em]
\addlinespace
5 & Ukuran jendela dan warna latar yang Anda pilih & \\[1.2em]
\addlinespace
6 & Posisi baris teks tambahan dan alasan pemilihannya & \\[1.2em]
\addlinespace
7 & Pesan penyimpanan yang Anda pakai & \\[1.2em]
\addlinespace
8 & Satu kendala yang Anda hadapi beserta cara penyelesaiannya & \\[2.0em]
\bottomrule
\end{tabularx}}
\end{table}

\begin{catatabox}
Praktikum ini dinilai sebagai artefak portofolio. Yang diperiksa adalah
kesesuaian fungsional program, kerapian kode, dan kelengkapan dokumentasi.
Rubrik penilaiannya ada pada Bagian~\ref{sec:asesmen-lingkungan}.
\end{catatabox}
